
\subsubsection{4.1 Research Questions}

\textbf{RQ1:} Is partial Spearman $\rho$ (MMLU-controlled) between BBQ-Disambig and BBQ-Ambig model rankings significantly positive (threshold: $\rho$ > 0.4, p < 0.05, N $\geq$ 10)?

\textbf{RQ2:} Does fairness rank stability exceed adversarial robustness rank stability by $\Delta\rho \geq$ 0.200 after capability control (Fisher z p < 0.05, N = 13)?

\textbf{RQ3:} Is partial $\rho$ for each adversarial robustness pair individually non-significant ($\rho$ < 0.4 or p $\geq$ 0.05), consistent with the adversarial disruption hypothesis?

RQ3 is presented before RQ2 in the Results section because the mechanism test result reframes the interpretation of the $\Delta\rho$ analysis: only after establishing that robustness rankings are also highly stable does the directional $\Delta\rho$ finding carry its correct evidential weight.

\subsubsection{4.2 Preregistered Criteria}

\textbf{P1 (Confirmatory, RQ1):} $\rho_{fairness}$ > 0.4 AND p < 0.05 AND N $\geq$ 10. Gate type: MUST\_WORK.

\textbf{P2 (Exploratory, RQ2):} $\Delta\rho \geq$ 0.200 AND Fisher z p < 0.05. Gate type: SHOULD\_WORK.

\textbf{P3 (Mechanism, RQ3):} $\rho_{robustness}$ pair not significantly positive ($\rho$ < 0.4 or p $\geq$ 0.05 for both pairs). Gate type: SHOULD\_WORK.

\subsubsection{4.3 Baselines}

Random permutation (expected $\rho$ = 0 under null; permutation test, 10,000 iterations) and raw (unadjusted) Spearman $\rho$ as comparison for quantifying the capability confound.

